\documentclass[../main.tex]{subfiles}
\begin{document}

\chapter{ReOrder Buffer}
\lessoninfo{
  video={Lesson 8, videos 1--40},
  textbook={H\&P \cite{hennessy2017} §3.6, §3.8, §C.4; \cite{smith1988}},
  keywords={exception, precise exception, phantom exception, reorder buffer, commit pointer, issue pointer, commit, branch misprediction recovery, register alias table, unified reservation station, superscalar, out-of-order execution},
}

Tomasulo's algorithm (Lesson 7) executes each instruction as soon as its
operands are available and writes its result into the register file when the
result is broadcast.  Renaming keeps the computed values correct, but two
events that every real processor must handle break the scheme: exceptions and
branch mispredictions.  In both cases instructions that follow the offending
instruction in program order may already have overwritten registers, and these
writes cannot be undone.  This lesson introduces the reorder buffer, which
lets instructions execute out of order but updates the registers strictly in
program order.  It traces the resulting processor in detail and ends with
unified reservation stations, superscalar execution, and a precise account of
which steps of an out-of-order processor actually work out of order.

\section{Exceptions in out-of-order execution}
\label{sec:l08-exceptions-ooo}

An exception is an event, such as a division by zero or an access to an
invalid memory address, that an instruction raises when it
executes.\source{Lesson 8, videos 1--2; H\&P §C.4, §3.6.}  The processor
saves the address of the excepting instruction and transfers control to an
exception handler, usually part of the operating system.  When the handler
has dealt with the cause, it may return to the saved address, so that the
excepting instruction is executed again and the program continues after it.
For the handler to work, and for the program to continue correctly after the
return, the registers and memory must look as they would in a processor that
executes one instruction at a time.

\begin{definition}[Precise exception]
  \label{def:l08-precise}
  An exception is a \term{precise exception}[正確な例外] if, when the
  handler starts, every instruction that precedes the excepting instruction
  in program order has completed and updated the registers and memory, and
  neither the excepting instruction nor any instruction that follows it has
  changed them.
\end{definition}

Tomasulo's algorithm does not provide precise exceptions.  Consider:
\begin{lstlisting}[language={[x86masm]Assembler}, numbers=none]
I1:  DIV R1, R2, R3     ; R3 = 0: division by zero
I2:  ADD R2, R4, R5     ; overwrites R2, read by I1
I3:  ADD R6, R6, R7     ; R6 <- R6 + R7
\end{lstlisting}
I2 and I3 do not depend on I1, so they execute and write R2 and R6 within a
few cycles of being issued.  The division takes many cycles, and in the model
used here the divider reports the division by zero when the operation
completes (\cref{fig:l08-exception}).  By then both later results are in the
register file, and three things go wrong.
\begin{itemize}
  \item The handler sees values of R2 and R6 that no sequential execution
    stopped at I1 could produce.
  \item After the return, I1 is executed again and now reads the value of R2
    written by I2, which in program order it must never see.
  \item I2 and I3 are also executed a second time.  I2 merely recomputes the
    same value, but I3 adds R7 to R6 twice.
\end{itemize}

\begin{figure}[htbp]
  \centering
  % Tomasulo's algorithm and an exception: two later instructions write their
% results into the register file long before the division reports the
% division by zero.
\begin{tikzpicture}[x=0.56cm,y=0.72cm, font=\footnotesize\sffamily, >={Stealth[length=1.8mm]},
  cell/.style={draw, fill=ln-box, minimum width=0.56cm, minimum height=0.44cm,
    inner sep=0pt, font=\scriptsize\sffamily},
  wr/.style={cell, draw=ln-caution, fill=ln-caution!15}]
  % cycle axis
  \foreach \c in {1,...,11} {
    \node[font=\scriptsize\sffamily, text=ln-muted] at (\c,0.68) {\c};
  }
  \node[anchor=east, font=\scriptsize\sffamily, text=ln-muted] at (0.4,0.68)
    {\iflangja{サイクル}{cycle}};
  % instruction labels
  \node[anchor=east] at (0.4,0)  {I1\enspace\texttt{DIV R1,R2,R3}};
  \node[anchor=east] at (0.4,-1) {I2\enspace\texttt{ADD R2,R4,R5}};
  \node[anchor=east] at (0.4,-2) {I3\enspace\texttt{ADD R6,R6,R7}};
  % I1: issue, long division, exception
  \node[cell] at (1,0) {I};
  \draw[fill=ln-box] (1.5,-0.22) rectangle (10.5,0.22);
  \node[font=\scriptsize\sffamily] at (6,0) {E\enspace(\iflangja{除算}{division})};
  \node[cell, draw=ln-caution, fill=ln-caution, text=white, font=\scriptsize\bfseries\sffamily]
    (exc) at (11,0) {!};
  % I2 and I3: short, independent
  \node[cell] at (2,-1) {I}; \node[cell] at (3,-1) {E}; \node[wr] (w2) at (4,-1) {W};
  \node[cell] at (3,-2) {I}; \node[cell] at (4,-2) {E}; \node[wr] (w3) at (5,-2) {W};
  % annotations
  \node[anchor=west, font=\scriptsize\sffamily, text=ln-caution] at (4.55,-1) {R2};
  \node[anchor=west, font=\scriptsize\sffamily, text=ln-caution] at (5.55,-2) {R6};
  \draw[ln-caution, thick, dashed] (11,-0.3) -- (11,-2.35);
  \node[anchor=north east, align=right, font=\scriptsize\sffamily, text=ln-caution]
    at (11.35,-2.35)
    {\iflangja{0 除算を報告：R2 と R6 はすでに上書き済み}%
              {division by zero reported: R2 and R6 already overwritten}};
\end{tikzpicture}

  \caption{An exception under Tomasulo's algorithm (I: issue, E: execute,
    W: write result).  I2 and I3 write R2 and R6 long before the division I1
    reports the division by zero.}
  \label{fig:l08-exception}
\end{figure}

\section{Branch mispredictions}
\label{sec:l08-mispredictions}

Branch mispredictions cause the same problem.\source{Lesson 8, video 3; H\&P
§3.6.}  In the following program the branch predictor lets the processor
fetch and issue I3 to I6 while I2 waits for the result of the division:
\begin{lstlisting}[language={[x86masm]Assembler}, numbers=none]
I1:        DIV R1, R2, R3     ; takes many cycles
I2:        BEQ R1, R4, Skip   ; predicted not taken
I3:        ADD R5, R6, R7     ; I3-I6: executed only
I4:        LW  R8, 0(R5)      ;   if not taken
I5:        MUL R9, R5, R5
I6:        SUB R6, R9, R7
           ...
     Skip: ...
\end{lstlisting}
I2 needs R1 and cannot execute before the division completes, whereas I3 has
its operands at once and I4 a cycle after I3.  Under Tomasulo's algorithm R5
and R8 are written long before the branch is resolved.  If the branch turns
out to be taken, the instructions of the wrong path can be removed from the
reservation stations and fetching can restart at \texttt{Skip}, but R5 and R8
have been overwritten with values that the program never computes, and their
correct values are lost.

A second problem arises when an instruction on the wrong path raises an
exception.  The program uses R5 as an address only when the branch is not
taken; when it is taken, the value $\text{R6}+\text{R7}$ need not be a valid
address, and the load I4 may raise an exception long before I2 executes.

\begin{definition}[Phantom exception]
  \label{def:l08-phantom}
  An exception raised by an instruction that lies on a mispredicted path, and
  therefore is never executed in a sequential execution of the program, is
  called a \term{phantom exception}[ファントム例外].
\end{definition}

A phantom exception must not reach the handler: in program order, the
instruction that raised it does not exist.

\section{Correct out-of-order execution}
\label{sec:l08-correct}

Both problems have the same cause.  An instruction changes the registers as
soon as it has computed its result, before it is known that every earlier
instruction has completed without an exception or a
misprediction.\source{Lesson 8, video 4; H\&P §3.6; \cite{smith1988}.}  The
remedy is to separate computing a result from making it visible:
\begin{itemize}
  \item instructions still \emph{execute} out of order, as soon as their
    operands are ready;
  \item results are still \emph{broadcast} out of order, so that waiting
    instructions receive them without delay;
  \item but results are \emph{written into the registers} in program order,
    and only once it is certain that the instruction belongs to the program's
    execution.
\end{itemize}
Holding the results until they may be written requires a structure that
remembers the program order of the instructions after they have been issued.
The reservation stations cannot do this: an instruction leaves its
reservation station as soon as its result is broadcast, and the stations are
allocated in no particular order.

\begin{keypoint}
  Execute out of order, broadcast out of order, but write the registers in
  program order.
\end{keypoint}

\section{The reorder buffer}
\label{sec:l08-rob}

The structure that remembers the program order is a queue of the issued
instructions.\source{Lesson 8, videos 5--10; H\&P §3.6.}

\begin{definition}[Reorder buffer]
  \label{def:l08-rob}
  The \term{reorder buffer}[リオーダバッファ]\index{ROB|see{reorder buffer}}
  (ROB) holds, in program order, every instruction that has been issued but
  has not yet passed the commit step of \cref{def:l08-commit}, in which its
  result is written into the register file.  Each entry has three fields:
  \begin{description}
    \item[Reg] the destination register of the instruction;
    \item[Value] its result, once computed;
    \item[Done] a bit that is set when Value holds the result.
  \end{description}
\end{definition}

The lecture draws only these three fields.  An entry must also record whether
the instruction raised an exception or is a mispredicted branch, and the
address of the instruction, which commit needs in order to take the exception
or to restart fetching (\cref{sec:l08-recovery,sec:l08-rob-exceptions}).

The ROB is a circular buffer managed by two pointers
(\cref{fig:l08-rob}).  The \term{issue pointer}[発行ポインタ] names the entry
that the next issued instruction receives, and it advances by one entry for
every issued instruction.  The oldest instruction that has not yet committed
is found at the \term{commit pointer}[コミットポインタ].  The entries from the
commit pointer up to, but not including, the issue pointer are in use and hold
instructions in program order; all others are free.  Since equal pointers
would describe both an empty and a full buffer, an occupancy count or a full
bit distinguishes the two.  In normal operation both pointers only advance,
wrapping around from the last entry to the first; recovery moves the issue
pointer back (\cref{sec:l08-recovery}).  When every entry is in use, no
instruction can be issued until the oldest one leaves the ROB.  Every issued
instruction receives an entry, including a branch, which writes no register:
its entry keeps its place in the program order and, as shown below, records
whether it was mispredicted.

\begin{figure}[htbp]
  \centering
  % The reorder buffer as a circular buffer: entries from the commit pointer
% up to the issue pointer are in use, in program order, and wrap around from
% the last entry to the first.
\begin{tikzpicture}[x=1cm,y=1cm, font=\footnotesize\sffamily, >={Stealth[length=1.8mm]},
  hd/.style={font=\footnotesize\bfseries\sffamily},
  used/.style={draw, fill=ln-accent!12, minimum height=0.45cm, inner sep=0pt},
  free/.style={draw, fill=white, minimum height=0.45cm, inner sep=0pt, text=ln-muted}]
  % column headings
  \node[hd] at (0.45,0.45) {Reg};
  \node[hd] at (1.55,0.45) {Value};
  \node[hd] at (2.65,0.45) {Done};
  % rows: entry / reg / value / done / used?
  \foreach \r/\reg/\val/\done/\st in {1/R2/{}/0/used, 2/R5/8/1/used,
      3/{}/{}/{}/free, 4/{}/{}/{}/free, 5/{}/{}/{}/free,
      6/R3/{}/0/used, 7/R1/17/1/used, 8/R3/25/1/used} {
    \pgfmathsetmacro{\y}{-(\r-1)*0.45}
    \node[\st, minimum width=0.9cm] at (0.45,\y) {\reg};
    \node[\st, minimum width=1.3cm] at (1.55,\y) {\val};
    \node[\st, minimum width=0.9cm] at (2.65,\y) {\done};
    \node[anchor=east, font=\scriptsize\sffamily, text=ln-muted] at (-0.05,\y) {ROB\r};
  }
  % pointers
  \draw[<-, thick, ln-accent] (3.15,-2.25) -- (4.0,-2.25)
    node[right, text=ln-accent] {\iflangja{コミットポインタ}{commit pointer}};
  \draw[<-, thick, ln-accent] (3.15,-0.9) -- (4.0,-0.9)
    node[right, text=ln-accent] {\iflangja{発行ポインタ}{issue pointer}};
  % wrap-around
  \draw[->, ln-muted] (-0.85,-3.15) -- (-1.15,-3.15) -- (-1.15,0.0) -- (-0.85,0.0);
  \node[rotate=90, font=\scriptsize\sffamily, text=ln-muted] at (-1.4,-1.57)
    {\iflangja{循環}{wrap around}};
  % program order of the entries in use
  \node[anchor=west, align=left, font=\scriptsize\sffamily, text=ln-muted] at (4.0,-3.15)
    {\iflangja{使用中のエントリ（プログラム順）：\\ROB6, ROB7, ROB8, ROB1, ROB2}%
              {entries in use, in program order:\\ROB6, ROB7, ROB8, ROB1, ROB2}};
\end{tikzpicture}

  \caption{A reorder buffer with eight entries.  The shaded entries, from the
    commit pointer up to the issue pointer, are in use; here this range wraps
    around from ROB8 to ROB1.}
  \label{fig:l08-rob}
\end{figure}

\subsection{Issue, dispatch and write result}

With a ROB, the three steps of Tomasulo's algorithm change as
follows.\index{Tomasulo's algorithm}
\begin{description}
  \item[Issue] The instruction needs a free reservation station (RS) and the
    free ROB entry at the issue pointer; if either is missing, issue waits.
    The ROB entry receives the destination register and a cleared Done bit,
    and the issue pointer advances.  Each operand is then looked up in the
    RAT, whose entries now name ROB entries instead of RSs:
    \begin{itemize}
      \item if the RAT entry points to the register file, the value is
        copied from the register;
      \item if it names a ROB entry whose Done bit is set, the value is copied
        from that entry;
      \item otherwise the RS records the ROB entry as the tag of the operand
        and waits for a broadcast with that tag.
    \end{itemize}
    Finally, the RAT entry of the destination register is set to the new ROB
    entry, and the RS records that entry as the tag of its result.
  \item[Dispatch] As before, an instruction whose operands are all available
    starts executing.  Its RS is freed at this point.
  \item[Write result] The result is broadcast on the CDB together with its
    tag, the ROB entry.  Waiting RSs capture it as before, and the ROB
    entry stores it and sets its Done bit.  Neither the register file nor the
    RAT is changed.
\end{description}

The second lookup case is new.  A result stays in the ROB from write result
until commit, and during that time the RAT keeps naming the ROB entry; an
instruction issued in that interval must find the value there, because the
broadcast it would otherwise wait for has already happened.

The earlier release of RSs is also new.\index{reservation station}  In
Tomasulo's algorithm an RS must stay allocated until write result, because
the station itself is the tag under which the result is broadcast.  With a ROB
the tag is the ROB entry, so the RS is no longer needed once the
operation has started, and the same number of RSs can hold more waiting
instructions.

\subsection{Commit}

A fourth step completes every instruction.

\begin{definition}[Commit]
  \label{def:l08-commit}
  In the \term{commit}[コミット]\index{retirement|see{commit}} step the
  processor examines the entry at the commit pointer.  If its Done bit is
  set, the value is written into the destination register, and if the RAT
  entry of that register still names this ROB entry, the RAT entry is reset
  to point to the register file.  The commit pointer then advances, which
  frees the entry.
\end{definition}

If the instruction at the commit pointer is not done, nothing commits, even if
later entries are: results enter the register file only in program order.  The
RAT entry is reset only when it still names the committing entry; otherwise a
later instruction writes the same register, and the RAT must keep naming that
instruction's entry (\cref{sec:l08-rat}).  Commit is also called
\emph{retirement}.

\subsection{Hardware organization}

\Cref{fig:l08-organization} shows the organization of Lesson 7 with a ROB
added.  The CDB now delivers results to the RSs and to the ROB, the RAT names
ROB entries, and the register file is written only by commit.  An instruction
passes through six steps: fetch, decode, issue, execute (from dispatch until
the operation completes), write result and commit.

\begin{figure}[htbp]
  \centering
  % Hardware organization of Tomasulo's algorithm with a reorder buffer: the RAT
% names ROB entries, the CDB delivers results to the RSs and the ROB, and only
% commit writes the register file.
\begin{tikzpicture}[x=1cm,y=1cm, font=\footnotesize\sffamily, >={Stealth[length=1.8mm]},
  row/.style={draw, fill=ln-box, minimum height=0.42cm, inner sep=0pt,
    font=\scriptsize\sffamily},
  lab/.style={font=\scriptsize\sffamily, text=ln-muted},
  hd/.style={font=\footnotesize\bfseries\sffamily},
  stepl/.style={font=\scriptsize\itshape\sffamily, text=ln-accent},
  unit/.style={draw, fill=white, trapezium, trapezium angle=65, shape border rotate=180,
    minimum width=1.3cm, minimum height=0.55cm, inner sep=1pt}]
  % ---------------------------------------------------------------- IQ
  \foreach \r in {1,...,4} {
    \node[row, minimum width=1.9cm] at (0.95,{-(\r-1)*0.42-0.21}) {};
  }
  \node[hd] at (0.95,0.25) {IQ};
  % ---------------------------------------------------------------- RAT
  \foreach \r/\v in {1/ROB2, 2/--, 3/ROB5, 4/ROB3} {
    \node[row, minimum width=0.95cm] at (3.45,{-(\r-1)*0.42-0.21}) {\v};
    \node[lab, anchor=east] at (2.95,{-(\r-1)*0.42-0.21}) {R\r};
  }
  \node[hd] at (3.45,0.25) {RAT};
  % ---------------------------------------------------------------- REGS
  \foreach \r/\v in {1/5, 2/-3, 3/12, 4/7} {
    \node[row, fill=white, minimum width=0.95cm] at (5.45,{-(\r-1)*0.42-0.21}) {\v};
    \node[lab, anchor=east] at (4.95,{-(\r-1)*0.42-0.21}) {R\r};
  }
  \node[hd] at (5.45,0.25) {\iflangja{レジスタ}{REGS}};
  % ---------------------------------------------------------------- ROB
  \node[lab] at (7.8,0.25) {Reg};
  \node[lab] at (8.8,0.25) {Value};
  \node[lab] at (9.8,0.25) {Done};
  \foreach \r/\reg/\val/\done/\f in {1/{}/{}/{}/white, 2/R1/9/1/ln-box, 3/R4/{}/0/ln-box,
      4/R3/{}/0/ln-box, 5/R3/21/1/ln-box, 6/{}/{}/{}/white} {
    \pgfmathsetmacro{\y}{-(\r-1)*0.42-0.21}
    \node[row, fill=\f, minimum width=0.8cm] at (7.8,\y) {\reg};
    \node[row, fill=\f, minimum width=1.2cm] at (8.8,\y) {\val};
    \node[row, fill=\f, minimum width=0.8cm] at (9.8,\y) {\done};
    \node[lab, anchor=west] at (10.25,\y) {ROB\r};
  }
  \node[hd, anchor=south] at (8.8,0.45) {ROB};
  \draw[->, thick, ln-accent] (7.38,-0.63) -- (6.65,-0.63) -- (6.65,-0.21) -- (5.97,-0.21);
  \node[stepl, anchor=north] at (6.65,-0.66) {\iflangja{コミット}{commit}};
  % ---------------------------------------------------------------- RSs
  \foreach \r in {1,2,3} {
    \node[row, minimum width=2.4cm] at (1.9,{-2.9-(\r-1)*0.42}) {};
    \node[lab, anchor=east] at (0.65,{-2.9-(\r-1)*0.42}) {AD\r};
  }
  \foreach \r in {1,2} {
    \node[row, minimum width=2.4cm] at (5.7,{-2.9-(\r-1)*0.42}) {};
    \node[lab, anchor=east] at (4.45,{-2.9-(\r-1)*0.42}) {MD\r};
  }
  % ---------------------------------------------------------------- units
  \node[unit] (add) at (1.9,-4.75) {$+,-$};
  \node[unit] (mul) at (5.7,-4.75) {$\times,\div$};
  \draw[->] (1.9,-3.95) -- (add.north);
  \draw[->] (5.7,-3.53) -- (mul.north);
  \node[stepl, anchor=west] at (1.95,-4.2) {\iflangja{ディスパッチ}{dispatch}};
  \node[stepl, anchor=west] at (5.75,-4.0) {\iflangja{ディスパッチ}{dispatch}};
  % ---------------------------------------------------------------- CDB
  \draw[line width=1.6pt, ln-accent] (1.9,-5.6) -- (8.8,-5.6);
  \draw[->] (add.south) -- (1.9,-5.6);
  \draw[->] (mul.south) -- (5.7,-5.6);
  \node[hd, text=ln-accent, anchor=north] at (4.0,-5.65) {CDB};
  \node[stepl, anchor=north west] at (5.8,-5.65) {\iflangja{結果書き込み}{write result}};
  % CDB to the RSs and to the ROB
  \draw[->, ln-accent] (3.55,-5.6) -- (3.55,-3.32) -- (3.12,-3.32);
  \draw[->, ln-accent] (7.35,-5.6) -- (7.35,-3.11) -- (6.92,-3.11);
  \draw[->, ln-accent] (8.8,-5.6) -- (8.8,-2.54);
  % ---------------------------------------------------------------- issue
  \draw[thick] (0.95,-1.7) -- (0.95,-2.31) -- (7.38,-2.31);
  \draw[->, thick] (7.0,-2.31) -- (7.38,-2.31);
  \draw[->, thick] (1.9,-2.31) -- (1.9,-2.68);
  \draw[->, thick] (5.7,-2.31) -- (5.7,-2.68);
  \draw[<->, dashed] (3.45,-1.7) -- (3.45,-2.31);
  \draw[->, dashed] (5.45,-1.7) -- (5.45,-2.31);
  \draw[->, dashed] (7.38,-1.89) -- (6.75,-1.89) -- (6.75,-2.31);
  \node[stepl, anchor=south east, inner xsep=1pt] at (7.36,-1.87)
    {\iflangja{Done なら値}{value if Done}};
  \node[stepl, anchor=south west] at (1.0,-2.31) {\iflangja{発行}{issue}};
\end{tikzpicture}

  \caption{Tomasulo's organization with a reorder buffer.  Issue allocates an
    RS and a ROB entry and renames the operands through the RAT; operand
    values come from the register file or from a Done ROB entry.  Results go
    from the CDB to the RSs and the ROB; commit copies the value at the commit
    pointer into the register file.}
  \label{fig:l08-organization}
\end{figure}

\begin{note}
  At any moment the most recent value of an architectural register, in the
  program order of the instructions issued so far, is in one of three places:
  in the register file, when the RAT points there; in a ROB entry whose Done
  bit is set; or not yet computed, in a ROB entry whose Done bit is clear.  The
  register file itself always holds the values as of the last committed
  instruction.
\end{note}

\section{Branch misprediction recovery}
\label{sec:l08-recovery}

With a ROB, instructions on a mispredicted path cannot change the register
file as long as they are removed before they commit.\source{Lesson 8, videos
11--12; H\&P §3.6.}  In the program of \cref{sec:l08-mispredictions}, I3 to
I6 are issued while I2 waits for R1, and all of them except I6, which is
still executing, write their results into their ROB entries.  When I2 finally executes and
the processor finds that the branch was mispredicted, the misprediction is
recorded in the ROB entry of I2, and nothing else happens yet.  Instructions
before the branch are unaffected and continue to execute and commit.

When the branch reaches the commit pointer, every instruction before it has
committed, so the register file holds exactly the values that a sequential
execution has after the branch.  Every instruction still in the processor
follows the branch and is on the wrong path.  The processor now performs
\term{branch misprediction recovery}[分岐予測ミス回復], which consists of four
actions (\cref{fig:l08-recovery}):
\begin{enumerate}
  \item discard all ROB entries by setting the issue pointer equal to the
    commit pointer, which leaves the ROB empty;
  \item reset every RAT entry to point to the register file;
  \item free all RSs and abandon the operations in the execution units, whose
    results could only go to discarded ROB entries;
  \item restart fetching at the correct successor of the branch.
\end{enumerate}
A correctly predicted branch simply commits: the commit pointer advances and
nothing else happens.

\begin{figure}[htbp]
  \centering
  % Branch misprediction recovery when the mispredicted branch reaches the commit
% pointer: every later ROB entry is on the wrong path and is discarded, and
% every RAT entry is reset to the register file.
\begin{tikzpicture}[x=1cm,y=1cm, font=\footnotesize\sffamily, >={Stealth[length=1.8mm]},
  row/.style={draw, fill=ln-box, minimum height=0.45cm, inner sep=0pt},
  lab/.style={font=\scriptsize\sffamily, text=ln-muted},
  hd/.style={font=\scriptsize\bfseries\sffamily},
  note/.style={font=\scriptsize\sffamily, anchor=west}]
  % ROB
  \node[hd] at (1.25,0.42) {\iflangja{命令}{Instruction}};
  \node[hd] at (2.9,0.42) {Done};
  \foreach \r/\ins/\done/\f in {1/{}/{}/white, 2/{BEQ R1,R4,Skip}/1/ln-box,
      3/{ADD R5,R6,R7}/1/ln-caution!12, 4/{LW~~R8,0(R5)}/1/ln-caution!12,
      5/{MUL R9,R5,R5}/1/ln-caution!12, 6/{SUB R6,R9,R7}/0/ln-caution!12,
      7/{}/{}/white, 8/{}/{}/white} {
    \pgfmathsetmacro{\y}{-(\r-1)*0.45}
    \node[row, fill=\f, minimum width=2.5cm] at (1.25,\y) {\texttt{\ins}};
    \node[row, fill=\f, minimum width=0.8cm] at (2.9,\y) {\done};
    \node[lab, anchor=east] at (-0.05,\y) {ROB\r};
  }
  \node[note, text=ln-accent] at (3.35,-0.45) {\iflangja{予測ミス}{mispredicted}};
  \node[note, text=ln-caution] at (3.35,-1.35) {\iflangja{例外}{exception}};
  \node[note, text=ln-muted] at (3.35,-2.25) {\iflangja{実行中}{executing}};
  % pointers
  \draw[->, thick, ln-accent] (-1.45,-0.45) -- (-0.9,-0.45);
  \node[font=\scriptsize\sffamily, text=ln-accent, anchor=east] at (-1.45,-0.45)
    {\iflangja{コミット}{commit}};
  \draw[->, thick, ln-accent] (-1.45,-2.7) -- (-0.9,-2.7);
  \node[font=\scriptsize\sffamily, text=ln-accent, anchor=east] at (-1.45,-2.7)
    {\iflangja{発行}{issue}};
  % wrong path
  \draw[thick, ln-caution] (4.75,-0.72) -- (4.9,-0.72) -- (4.9,-2.43) -- (4.75,-2.43);
  \node[font=\scriptsize\sffamily, text=ln-caution, align=left, anchor=west] at (4.95,-1.57)
    {\iflangja{誤った経路：\\すべて破棄}{wrong path:\\all discarded}};
  % RAT before and after
  \node[hd] at (7.7,0.42) {\iflangja{回復前}{before}};
  \node[hd] at (8.7,0.42) {\iflangja{回復後}{after}};
  \node[hd, anchor=south] at (8.2,0.62) {RAT};
  \foreach \r/\reg/\before in {1/R1/--, 2/R5/ROB3, 3/R6/ROB6, 4/R8/ROB4, 5/R9/ROB5} {
    \pgfmathsetmacro{\y}{-(\r-1)*0.45}
    \node[lab, anchor=east] at (7.15,\y) {\reg};
    \node[row, minimum width=1.0cm, font=\scriptsize\sffamily] at (7.7,\y) {\before};
    \node[row, fill=white, minimum width=1.0cm, font=\scriptsize\sffamily] at (8.7,\y) {--};
  }
  \node[lab, align=left, anchor=north west] at (6.8,-2.2)
    {\iflangja{--：レジスタ\\\phantom{--：}ファイルを指す}{--: points to the\\\phantom{--: }register file}};
\end{tikzpicture}

  \caption{Recovery when the mispredicted branch I2 reaches the commit pointer.
    I1 has already committed; the ROB entries of I3 to I6 are discarded,
    including the exception recorded for I4, and all RAT entries return to
    the register file.}
  \label{fig:l08-recovery}
\end{figure}

\begin{note}
  Recovering at commit is the simplest scheme and the one used in the lecture.
  Real processors start recovery as soon as the mispredicted branch executes,
  discarding only the ROB entries after the branch and restarting fetching at
  once, so that no cycles are lost while the branch waits to commit (H\&P
  §3.6).  The RAT must then be brought back to its state at the branch, for
  example from a copy saved when the branch was issued.
\end{note}

\section{Exceptions with a reorder buffer}
\label{sec:l08-rob-exceptions}

Exceptions are handled by treating an exception as one more result of the
instruction that raises it.\source{Lesson 8, videos 13--15; H\&P §3.6.}  When
an instruction raises an exception during execution, the processor records
the exception in the instruction's ROB entry instead of invoking the handler.
When that entry reaches the commit pointer, the exception is taken: the
processor discards the ROB entries, resets the RAT, and frees the RSs and
execution units exactly as in misprediction recovery, saves the address of the
excepting instruction, and starts fetching the exception handler.

This removes both problems of
\cref{sec:l08-exceptions-ooo,sec:l08-mispredictions}.
\begin{description}
  \item[Precise exceptions] When the excepting instruction reaches the commit
    pointer, all earlier instructions have committed and no later instruction
    has written the register file.  Stores likewise write memory only when
    they commit (Lesson~9), so memory is also unchanged by later instructions.
    The handler sees exactly the state of \cref{def:l08-precise}; after the
    return, the later instructions are executed from that state.
  \item[No phantom exceptions] An exception on a mispredicted path is recorded
    in an entry that follows the mispredicted branch.  The branch commits
    first, and its recovery discards the entry together with the exception,
    which is never taken.  In \cref{fig:l08-recovery} the exception of I4
    disappears when I2 commits.
\end{description}

\subsection{The outside view of execution}
\label{sec:l08-outside}

The ROB separates two notions of when an instruction is executed.  Inside the
processor, an instruction executes when an execution unit performs its
operation, in any order.  The programmer, the operating system and the
exception handler, however, see only the register file, memory and the
program counter, and for them an instruction is executed when it commits.
Commits happen one instruction after another in program order, exactly as in a
processor that completes each instruction before starting the next.  An
instruction whose result is in the ROB but has not committed has, seen from
outside, not been executed yet; if its result is discarded, it has never been
executed at all.  The example of \cref{sec:l08-example} ends with an
instance.

\section{RAT updates on commit}
\label{sec:l08-rat}

The conditional RAT update of \cref{def:l08-commit} is needed because a
register may have several pending writers.\source{Lesson 8, video 16; H\&P
§3.6.}\index{register alias table}  Suppose the instructions in ROB2 and ROB5
both write R4, and ROB2 commits while ROB5 is still in the ROB.  The RAT entry
of R4 names ROB5, because that instruction was issued later.
\begin{itemize}
  \item Commit writes the value of ROB2 into R4.  This is the correct
    architectural value after ROB2, and it is the value an exception handler
    must see if an instruction between ROB2 and ROB5 raises an exception.
  \item The RAT entry must nevertheless keep naming ROB5.  An instruction
    issued now follows ROB5 in program order and must use the value of ROB5,
    not the older value that has just entered the register file.
\end{itemize}
When ROB5 commits, the RAT entry of R4 still names ROB5 unless an even later
writer has been issued in the meantime, and only then is it reset.  The
opposite mistake, leaving a RAT entry that names a committed entry, is equally
wrong: the entry is freed at commit and may soon be given to an unrelated
instruction.

\section{A reorder buffer example}
\label{sec:l08-example}

We now trace a complete program through a processor with a
ROB.\source{Lesson 8, videos 17--30; H\&P §3.6.}  The processor has three
RSs for addition and subtraction (AD1 to AD3), two for multiplication and
division (MD1, MD2) and a ROB of eight entries.  ADD and SUB take 1 cycle, MUL
4 cycles and DIV 6 cycles, and there are enough execution units that no ready
instruction waits for one.  The steps follow these conventions:
\begin{itemize}
  \item one instruction is issued per cycle, into the free RS of the required
    kind with the lowest number;
  \item an RS freed by a dispatch in cycle $C$, and a ROB entry freed by a
    commit in cycle $C$, can be given to a new instruction from cycle $C+1$;
  \item an instruction is dispatched no earlier than the cycle after its issue
    and the cycle after the broadcast of each operand it waits for;
  \item an operation dispatched in cycle $D$ with an execution time of $k$
    cycles executes in cycles $D$ to $D+k-1$ and writes its result in cycle
    $D+k$; there is one CDB, and in this example no two results are written
    in the same cycle;
  \item an instruction commits no earlier than the cycle after its write
    result, and at most one instruction commits per cycle.
\end{itemize}

\needspace{12\baselineskip}
\begin{example}
  \label{ex:l08-trace}
  The program is
  \begin{lstlisting}[language={[x86masm]Assembler}, numbers=none]
I1:  DIV R1, R2, R3     ; R1 <- R2 / R3
I2:  ADD R4, R5, R6     ; R4 <- R5 + R6
I3:  MUL R2, R4, R4     ; R2 <- R4 * R4
I4:  SUB R5, R1, R6     ; R5 <- R1 - R6
I5:  ADD R4, R4, R3     ; R4 <- R4 + R3
I6:  ADD R6, R2, R4     ; R6 <- R2 + R4
\end{lstlisting}
  with initial values $\text{R1}=5$, $\text{R2}=48$, $\text{R3}=4$,
  $\text{R4}=2$, $\text{R5}=7$ and $\text{R6}=3$.  All RAT entries point to
  the register file, and both ROB pointers are at ROB1.
\end{example}

\paragraph{Cycles 1--6: issue.}
One instruction is issued per cycle, I$k$ into ROB$k$.  I1 copies 48 and 4
from the register file into MD1 and sets the RAT entry of R1 to ROB1.  In
cycle~2 I1 is dispatched and frees MD1, and I2 is issued into AD1 with the
values 7 and~3.  I3, issued in cycle~3, finds ROB2 in the RAT entry of R4; ROB2
is not done, so both operands of I3 wait for ROB2.  I3 can use MD1 again,
which under Tomasulo's algorithm alone would stay occupied by I1 until
cycle~8.  I4 is issued into AD1, freed by I2 in cycle~3; its first operand
waits for ROB1, its second is 3.  In cycle~4 I2 writes its result: ROB2
receives 10, and I3 captures it.  In cycle~5 I5 is issued into AD2 and finds
the RAT entry of R4 still naming ROB2, which is done, so it copies 10 from the
ROB, not from the register file, which still holds the old value~2; I5 then
renames R4 to ROB5.  In
cycle~6 I6 is issued into AD3, because AD2, freed by the dispatch of I5 in the
same cycle, can be used only from cycle~7; I6 waits for ROB3 and ROB5.
\Cref{tab:l08-rob-state,tab:l08-rat-state} show the state at the end of
cycle~6: AD1 holds I4 and AD3 holds I6, all other RSs are free, and the
register file is unchanged.

\begin{table}[htbp]
  \centering\small
  \caption{ROB of \cref{ex:l08-trace} at the end of cycles 6 and 10.  At the
    end of cycle~6 the commit pointer is at ROB1, at the end of cycle~10 at
    ROB3; the issue pointer is at ROB7 in both.}
  \label{tab:l08-rob-state}
  \begin{tabular}{@{}llcrcrc@{}}
    \toprule
    & & & \multicolumn{2}{c}{Cycle 6} & \multicolumn{2}{c}{Cycle 10} \\
    \cmidrule(lr){4-5}\cmidrule(l){6-7}
    Entry & Instruction & Reg & Value & Done & Value & Done \\
    \midrule
    ROB1 & \texttt{DIV R1, R2, R3} & R1 &     & 0 & \multicolumn{2}{c}{committed} \\
    ROB2 & \texttt{ADD R4, R5, R6} & R4 & 10  & 1 & \multicolumn{2}{c}{committed} \\
    ROB3 & \texttt{MUL R2, R4, R4} & R2 &     & 0 & 100 & 1 \\
    ROB4 & \texttt{SUB R5, R1, R6} & R5 &     & 0 & 9   & 1 \\
    ROB5 & \texttt{ADD R4, R4, R3} & R4 &     & 0 & 14  & 1 \\
    ROB6 & \texttt{ADD R6, R2, R4} & R6 &     & 0 &     & 0 \\
    \bottomrule
  \end{tabular}
\end{table}

\begin{table}[htbp]
  \centering\small
  \caption{RAT and register file of \cref{ex:l08-trace}.  A dash means that
    the RAT entry points to the register file.}
  \label{tab:l08-rat-state}
  \begin{tabular}{@{}lcrcrr@{}}
    \toprule
    & \multicolumn{2}{c}{Cycle 6} & \multicolumn{2}{c}{Cycle 10} & End \\
    \cmidrule(lr){2-3}\cmidrule(lr){4-5}\cmidrule(l){6-6}
    Register & RAT & Value & RAT & Value & Value \\
    \midrule
    R1 & ROB1 & 5  & --   & 12 & 12  \\
    R2 & ROB3 & 48 & ROB3 & 48 & 100 \\
    R3 & --   & 4  & --   & 4  & 4   \\
    R4 & ROB5 & 2  & ROB5 & 10 & 14  \\
    R5 & ROB4 & 7  & ROB4 & 7  & 9   \\
    R6 & ROB6 & 3  & ROB6 & 3  & 114 \\
    \bottomrule
  \end{tabular}
\end{table}

\paragraph{Cycles 7--8: results, but no commit.}
I5 writes 14 into ROB5 in cycle~7 and I1 writes 12 into ROB1 in cycle~8; I6
and I4 capture these values.  Although ROB2 has been done since cycle~4 and
ROB5 since cycle~7, nothing commits, because ROB1, at the commit pointer, is
not done until the end of cycle~8.

\paragraph{Cycles 9--10: the first commits.}
In cycle~9 I1 commits: R1 becomes 12 and, since the RAT entry of R1 still names
ROB1, it is reset to the register file.  In the same cycle I3 writes 100 into
ROB3 and I4 is dispatched.  In cycle~10 I2 commits and R4 becomes 10, but the
RAT entry of R4 names ROB5 and is left unchanged, as \cref{sec:l08-rat}
requires.  I4 writes 9 and I6 is dispatched.

\paragraph{Cycles 11--14: the remaining commits.}
I6 writes 114 in cycle~11.  I3, I4, I5 and I6 commit in cycles 11 to 14, one
per cycle.  Each is the last writer of its register, so each resets its RAT
entry.  At the end the ROB is empty, all RAT entries point to the register
file, and the registers hold $\text{R1}=12$, $\text{R2}=100$, $\text{R3}=4$,
$\text{R4}=14$, $\text{R5}=9$ and $\text{R6}=114$, the values of a sequential
execution.

At the end of cycle~10 the register file holds 10 in R4 while the RAT sends
readers of R4 to ROB5, which holds 14.  Both are right: 10 is the value of R4
after I2, the last committed instruction, and 14 is the value after I5, the
last issued writer.

The example also illustrates the outside view of \cref{sec:l08-outside}.
Suppose that I4 raises an exception when it executes in cycle~9.  The exception
is recorded in ROB4, I2 and I3 still commit in cycles 10 and~11, and in
cycle~12 I4 reaches the commit pointer and the exception is taken.  The
handler finds $\text{R1}=12$, $\text{R2}=100$ and $\text{R4}=10$, but not
$\text{R4}=14$, although I5 wrote 14 into ROB5 in cycle~7, before I1 had
even finished: from outside, execution stopped exactly before I4.

\section{Timing with a reorder buffer}
\label{sec:l08-timing}

The cycle-by-cycle state of a ROB processor can be summarized by the cycle in
which each instruction passes each step, as for Tomasulo's algorithm in
Lesson~7, with one additional column for commit.\source{Lesson 8, videos
31--34; H\&P §3.6.}  We keep the notation of Lesson~7: $I_n$, $D_n$ and $W_n$
are the cycles of issue, dispatch and write result of instruction $n$, $k_n$
is its execution time, and $P(n)$ is the set of instructions that produce its
operands; $C_n$ denotes its commit cycle.  Under the conventions of
\cref{sec:l08-example}, $I_n$ is the first cycle after $I_{n-1}$ in which an
RS of the required kind and a ROB entry are free.  An RS released by the
dispatch of instruction $m$ is free from cycle $D_m+1$, and a ROB entry
released by the commit of $m$ from cycle $C_m+1$.  With enough execution units
and no competition for the CDB, the other cycles follow from
\begin{equation}
  D_n = 1 + \max\Bigl(I_n,\; \max_{p \in P(n)} W_p\Bigr), \qquad
  W_n = D_n + k_n, \qquad
  C_n = \max\bigl(W_n + 1,\; C_{n-1} + 1\bigr),
  \label{eq:l08-timing}
\end{equation}
where the first equation is the bound of \cref{eq:l07-dispatch} with a free
execution unit.\tip{Fill in the table row by row in program order, each row
from left to right, with Commit last.}  Unlike execution, commit depends on
the preceding instruction even when the two share no data.  The point at which
RSs are released and the number of commits per cycle are design assumptions
that change the table.  If an RS is released only at write result, as in
Lesson~7, it is free from cycle $W_m+1$; if two instructions may commit in the
same cycle,
$C_n = \max(W_n+1,\; C_{n-1},\; C_{n-2}+1)$.

\begin{table}[htbp]
  \centering\small
  \caption{Timing of \cref{ex:l08-trace}.  The Execute column gives the cycles
    $D_n$ to $D_n+k_n-1$ in which the operation is performed.}
  \label{tab:l08-timing}
  \begin{tabular}{@{}llcccc@{}}
    \toprule
    & Instruction & Issue & Execute & Write result & Commit \\
    \midrule
    I1 & \texttt{DIV R1, R2, R3} & 1 & 2--7 & 8  & 9  \\
    I2 & \texttt{ADD R4, R5, R6} & 2 & 3    & 4  & 10 \\
    I3 & \texttt{MUL R2, R4, R4} & 3 & 5--8 & 9  & 11 \\
    I4 & \texttt{SUB R5, R1, R6} & 4 & 9    & 10 & 12 \\
    I5 & \texttt{ADD R4, R4, R3} & 5 & 6    & 7  & 13 \\
    I6 & \texttt{ADD R6, R2, R4} & 6 & 10   & 11 & 14 \\
    \bottomrule
  \end{tabular}
\end{table}

\Cref{tab:l08-timing} lists the result for the example.  The Execute and
Write result columns are out of order: I5 writes its result before I1, I3 and
I4, which precede it.  The Commit column is increasing.  I2 and I5 commit six
cycles after writing their results, held up by the slow division at the head
of the ROB and, after it commits in cycle~9, by the limit of one commit per
cycle: every later instruction commits exactly one cycle after its
predecessor.  This waiting delays no computation: I3 and I5 obtained the
result of I2, and I6 those of I3 and I5, without waiting for any commit.  The
ROB delays only the moment at which results become part of the architectural
state.

\needspace{4\baselineskip}
The ROB can nevertheless limit execution when it is too small.  With only four
entries, I5 could not be issued before cycle~10, the cycle after I1 commits,
and it would execute several cycles later than in \cref{tab:l08-timing} although its
operands are available from cycle~5.  A ROB must therefore be large enough to
hold all instructions issued while a slow instruction waits to commit.

\section{Unified reservation stations}
\label{sec:l08-unified}

In \cref{fig:l08-organization} each kind of execution unit has its own group of
RSs, and the size of each group is fixed at design time.\source{Lesson 8,
video 35; H\&P ch.~3.}  Programs differ in their mix of operations, so a
program with many additions can find all adder RSs occupied, and issue must
wait, while the multiplier RSs are empty.  Because instructions are issued in
program order, even a multiplication that follows the waiting addition cannot
use the free multiplier RSs.

\begin{definition}[Unified reservation stations]
  \label{def:l08-unified}
  A \term{unified reservation station}[統合リザベーションステーション] is a
  single pool of RSs shared by all operations: an instruction can be issued
  into any free RS of the pool.
\end{definition}

Issue then waits only when all RSs are occupied, so the same number of RSs is
used more effectively.  The cost lies in dispatch: instead of choosing among
the RSs of one group for one execution unit, the dispatch logic must find, for
each free execution unit, the ready instructions in the whole pool whose
operation that unit performs, and select among them.

\section{Superscalar processors}
\label{sec:l08-superscalar}

Lesson~6 showed that the IPC of a processor is limited by the number of
instructions it can handle per cycle (\cref{sec:l06-ipc}).\source{Lesson 8,
video 36; H\&P §3.8.}

\begin{definition}[Superscalar processor]
  \label{def:l08-superscalar}
  A \term{superscalar}[スーパースカラ] processor performs every step on more
  than one instruction per cycle: it fetches, decodes and issues several
  instructions, dispatches several to execution units, writes several results
  and commits several instructions in each cycle.
\end{definition}

Every step needs its own hardware for this.  Fetch and decode must handle
several instructions at once; issue must rename several instructions in the
same cycle, each seeing the RAT updates of the ones before it; dispatch needs
several execution units; writing several results needs several CDBs, and
every RS operand must compare its tag against all of them; commit must
examine several entries at the commit pointer and write several registers.

The steps form a chain through which every instruction passes, so the
processor sustains no more instructions per cycle than its most restrictive
step handles.  A processor that fetches, decodes and dispatches four
instructions per cycle but issues three and commits two completes, in the
long run, at most two instructions per cycle, however much ILP the program
has.  Widening one step pays off only if that step is the bottleneck.

\section{Terminology}
\label{sec:l08-terminology}

The names used in this course for the steps of an out-of-order processor are
not universal.\source{Lesson 8, video 37; H\&P ch.~3.}  Papers, textbooks
and processor documentation use several alternatives
(\cref{tab:l08-terminology}).\caution{The same word can name different steps:
in some texts \emph{dispatch} is what this course calls issue, and
\emph{issue} what it calls dispatch.}  Some of them call commit
\emph{complete}; in this chapter, by contrast, an operation completes when its
execution ends.

\begin{table}[htbp]
  \centering\small
  \caption{Names of the steps in this course and in other texts.}
  \label{tab:l08-terminology}
  \begin{tabular}{@{}ll@{}}
    \toprule
    This course & Names in other texts \\
    \midrule
    Issue & issue, allocate, dispatch \\
    Dispatch & dispatch, execute, issue \\
    Commit & commit, complete, retire, graduate \\
    \bottomrule
  \end{tabular}
\end{table}

\section{What is out of order}
\label{sec:l08-order}

A processor with out-of-order execution\index{out-of-order execution} does not
perform all of its steps out of order (\cref{fig:l08-order}).\source{Lesson 8,
videos 38--40; H\&P §3.6.}
\begin{description}
  \item[Fetch and decode] take instructions in program order, along the
    predicted path.
  \item[Issue] is in program order.  Renaming requires it, because each
    instruction must see the RAT as left by all instructions before it, and
    ROB entries are allocated in program order.
  \item[Execute and write result] are out of order: an instruction starts as
    soon as its operands are ready and finishes when its operation completes.
    These are the only steps in which a later instruction overtakes an
    earlier one.
  \item[Commit] is in program order again; the ROB restores the order that
    execution relaxed.
\end{description}

\begin{figure}[htbp]
  \centering
  % The steps of an out-of-order processor: only execution and write result
% are performed out of program order.
\begin{tikzpicture}[x=1cm,y=1cm, font=\scriptsize\sffamily, >={Stealth[length=1.6mm]},
  stp/.style={draw, fill=ln-box, minimum width=1.55cm, minimum height=0.62cm,
    inner sep=1pt, align=center},
  ooo/.style={stp, draw=ln-accent, fill=ln-accent!15}]
  \foreach \i/\txt/\sty in {0/{\iflangja{フェッチ}{Fetch}}/stp,
      1/{\iflangja{デコード}{Decode}}/stp, 2/{\iflangja{発行}{Issue}}/stp,
      3/{\iflangja{実行}{Execute}}/ooo, 4/{\iflangja{結果書き込み}{Write result}}/ooo,
      5/{\iflangja{コミット}{Commit}}/stp} {
    \node[\sty] (s\i) at ({\i*1.95},0) {\txt};
  }
  \foreach \i [evaluate=\i as \j using int(\i+1)] in {0,...,4} {
    \draw[->] (s\i.east) -- (s\j.west);
  }
  \draw[thick] ([yshift=4pt]s0.north west) -- ++(0,4pt) -| ([yshift=4pt]s2.north east);
  \draw[thick, ln-accent] ([yshift=4pt]s3.north west) -- ++(0,4pt) -| ([yshift=4pt]s4.north east);
  \draw[thick] ([yshift=4pt]s5.north west) -- ++(0,4pt) -| ([yshift=4pt]s5.north east);
  \node[anchor=base] at ([yshift=11pt]s1.north) {\iflangja{プログラム順}{in program order}};
  \node[anchor=base, text=ln-accent] at ([yshift=11pt]$(s3.north)!0.5!(s4.north)$)
    {\iflangja{アウトオブオーダ}{out of order}};
  \node[anchor=base] at ([yshift=11pt]s5.north) {\iflangja{プログラム順}{in program order}};
\end{tikzpicture}

  \caption{The steps of an out-of-order processor with a ROB.  Only execution
    and write result take place out of program order.}
  \label{fig:l08-order}
\end{figure}

\begin{keypoint}[Lesson summary]
  Tomasulo's algorithm writes results into the registers as soon as they are
  broadcast, so exceptions are not precise, mispredicted instructions corrupt
  registers, and wrong-path instructions can raise phantom exceptions.  The
  reorder buffer keeps issued instructions in program order between a commit
  pointer and an issue pointer.  The RAT names ROB entries, results are
  broadcast to the RSs and the ROB, RSs are freed at dispatch, and only commit,
  one ROB entry after another in program order, writes the register file and
  resets a RAT entry that still names the committing entry.  A mispredicted
  branch or an excepting instruction is acted on when it commits: the ROB is
  emptied, the RAT points back to the register file, the RSs and execution
  units are cleared, and fetching restarts at the correct address.  Timing
  tables gain a commit column that follows \cref{eq:l08-timing}.  Unified
  reservation stations share one RS pool among all operations; a superscalar
  processor handles several instructions per cycle in every step and is limited
  by its narrowest step.  Only execution and write result are out of order.
\end{keypoint}

\end{document}
